% intro.tex

%%%%%%%%%%%%%%%
\begin{frame}{}
  \begin{definition}[Equivalence Relation]
    $R \subseteq X \times X$ is an \red{\it equivalence relation} on $X$ iff $R$ is
    \begin{itemize}
      \item reflexive (自反): $\forall a \in X\; a R a$
      \item symmetric (对称): $\forall a, b \in X\; (a R b \leftrightarrow b R a)$
      \item transitive (传递): $\forall a, b, c \in X\; (a R b \land b R c \to a R c)$
    \end{itemize}
  \end{definition}

	\pause
  \vspace{1em}
  \fig{width = 0.40\textwidth}{figs/partition}
  \[
    \red{\text{Equivalence Relation} \iff \text{Partition} \iff \text{Abstraction}}
  \]
\end{frame}
%%%%%%%%%%%%%%%

%%%%%%%%%%%%%%%
\begin{frame}{}
  \begin{definition}[Equivalence Relation]
    $R \subseteq X \times X$ is an {\it equivalence relation} on $X$ iff $R$ is
    \begin{itemize}
      \item reflexive (自反): $\forall a \in X\; a R a$
      \item symmetric (对称): $\forall a, b \in X\; (a R b \leftrightarrow b R a)$
      \item transitive (传递): $\forall a, b, c \in X\; (a R b \land b R c \to a R c)$
    \end{itemize}
  \end{definition}

	\begin{columns}
		\column{0.50\textwidth}
			\fig{width = 0.60\textwidth}{figs/cat-question-mark}
		\column{0.50\textwidth}
			\begin{center}
				\fig{width = 0.45\textwidth}{figs/ordering}
        \vspace{-1em}
				\red{Ordering?}
			\end{center}
	\end{columns}
\end{frame}
%%%%%%%%%%%%%%%

%%%%%%%%%%%%%%%
\begin{frame}{}
	\begin{quote}
		``\red{Order theory} is a branch of mathematics that
		investigates the \purple{intuitive notion of order}
		using \blue{binary relations}.
  \end{quote}

  \begin{columns}
    \column{0.50\textwidth}
      \fig{width = 0.40\textwidth}{figs/ordering}
    \column{0.50\textwidth}
      \uncover<3->{
        \fig{width = 0.80\textwidth}{figs/questions}
      }
  \end{columns}

  \pause
  \vspace{0.30cm}
  \begin{quote}
		It provides a formal framework for describing statements
		such as \teal{``this is less than that''}
    or \teal{``this precedes that''}.''

    \vspace{1em}
		\hfill --- \href{https://en.wikipedia.org/wiki/Order\_theory}{Order theory (Wikipedia)}
	\end{quote}
\end{frame}
%%%%%%%%%%%%%%%

%%%%%%%%%%%%%%%
\begin{frame}{}
  \fig{width = 0.70\textwidth}{figs/ordering-relations}
\end{frame}
%%%%%%%%%%%%%%%